\documentclass[11pt]{letter}
\usepackage[utf8]{inputenc}
\usepackage{geometry}
\usepackage{hyperref}
\geometry{a4paper, margin=1in}

\address{
  \textbf{Abdul Fatah} \\
  Department of Computer Science and Applied Physics \\
  Atlantic Technological University \\
  Galway, Ireland \\
  Email: \href{mailto:abdul.fatah@research.atu.ie}{abdul.fatah@research.atu.ie}
}

\signature{
  Abdul Fatah \\
  \textit{On behalf of the authors:} \\
  Abdul Fatah, Ian McLoughlin, and Saim Ghafoor \\
  Atlantic Technological University, Ireland
}

\begin{document}

\begin{letter}{
  Guest Editors \\
  ACM Transactions on Mathematical Software (TOMS) \\
  Special Issue: Mixed-Precisions Techniques for HPC Algorithms and Applications using Hardware Accelerators
}

\opening{Dear Guest Editors,}

We are pleased to submit our manuscript entitled \textbf{``A Reproducible Precision-Aware Workflow for Unanchored Approximate Mutually Unbiased Bases in Dimension Six''} for consideration for publication in the \textit{ACM Transactions on Mathematical Software (TOMS)} Special Issue on \textit{Mixed-Precisions Techniques for HPC Algorithms and Applications using Hardware Accelerators}.

Our work addresses a classical, longstanding challenge in quantum information theory and finite-dimensional Hilbert space geometry: the existence of a complete set of seven Mutually Unbiased Bases (MUBs) in dimension six ($d=6$). We approach this from a mathematical software perspective by introducing an unanchored optimization workflow that utilizes automatic differentiation and Lie-algebra exponentiation to search for approximate configurations.

We believe our manuscript is a perfect fit for this Special Issue, as it directly intersects mathematical software design, hardware acceleration, and precision-aware computing:

\begin{itemize}
    \item \textbf{Custom Hardware-Accelerated Exponentiation Layer:} To support GPU acceleration natively on both NVIDIA (CUDA) and Apple Silicon (MPS) backends, we developed a highly portable complex matrix Taylor-series exponentiation layer. This layer bypasses backend device limitations for complex matrix exponentials while maintaining truncation errors strictly below floating-point precision limits.
    \item \textbf{Precision-Aware Landscape Analysis:} We conduct a thorough analysis contrasting double-precision (\texttt{complex128}) reference runs with single-precision (\texttt{complex64}) implementations. We demonstrate that while qualitative structural transitions remain stable, the fine-grained classification of near-exact pairs, loss convergence metrics, and the optimization basin selection are highly sensitive to numerical floating-point precision, with single-precision execution prone to trapping in specific low-loss basins.
    \item \textbf{Multi-Dimensional Benchmarks and GPU Crossover:} We present extensive scaling benchmark campaigns comparing CPU and GPU runtimes across dimensions. We identify and empirically document a clear execution speed crossover point around dimension ninety ($d=90$).
    \item \textbf{Reproducible Software Workflow:} Adhering to the core philosophy of TOMS, our methodology is packaged as a modern, reproducible workflow leveraging Python, PyTorch, and \textit{marimo} interactive notebooks. All optimization seeds, overlap tensors, diagnostics, and resulting basis matrices are fully archived and publicly exportable, allowing independent verification and extension.
\end{itemize}

This manuscript has not been published elsewhere and is not under consideration by any other journal. All authors have read and approved the final manuscript for submission. 

Thank you very much for your time and for considering our work. We look forward to receiving the reviews.

\closing{Sincerely yours,}

\end{letter}
\end{document}
